% SPDX-License-Identifier: LPPL-1.3c
% Copyright (C) 2026 Christophe Jorssen
% Author and Current Maintainer: Christophe Jorssen <christophe.jorssen@gmail.com>
% LPPL maintenance status: maintained
% This file is part of luafemm. See LICENSE and LICENSES.md for its terms.
\section{Controlling mesh quality}
\label{sec:mesh-quality}
A smaller |mesh size| puts more points into the model, but it does not
necessarily improve the shape of every triangle. A narrow air gap or the
transition from a fine material region to coarse air can leave long, thin
elements. The optional |minimum angle| and |maximum area| keys address
triangle shape and area separately from nominal point spacing.

\subsection{Draft first, request quality second}
Keep both quality targets at zero while positioning parts and checking their
material declarations. Inspect the mesh, then add a reusable style with a
modest angle target, usually 20 degrees. Declare that style in the picture
options, together with the domain and spacing: refinement is part of meshing,
so it must be configured before the first mesh or field request.

The complete example below uses the same U electromagnet and 8 mm draft
spacing twice. Only the angle target changes. The two |femm/problem| keys
provided by the styles configure one model per picture; their options compose
before the model is created. |\pic {femm mesh};| performs meshing and quality
refinement, but does not solve the magnetic field.

For these sources, the minimum angle increases from about $7.29^\circ$ to
$20.00^\circ$ with 148 added points. The diagnostics are computed by the
example, rather than written into the figure. Counts are observations for
this geometry and revision, not promises for other devices.
\luafemmexample{mesh-quality}

\subsection{How the refinement proceeds}
The initial mesh still follows global and local point spacing. When either
quality target is nonzero, refinement first examines constrained segments.
A neighbouring triangle's opposite vertex encroaches on a segment if it lies
strictly inside the circle whose diameter is that segment. Such a segment
is bisected before processing deficient triangles. Points exactly on the
circle do not trigger subdivision; the disk test uses a filtered dot product
with exact-expansion fallback.

Deficient triangles are processed in order of shortest edge, with deterministic
tie breaking. Their proposed circumcentres are checked against the constrained
segments. If a candidate encroaches a segment, that segment is split instead.
The point-location walk also stops at a constrained interface and requests
its subdivision before continuing. Accepted insertions are followed by local
Lawson flips that preserve every constrained edge. Newly created triangles
and neighbouring segments are checked again. All split segments retain their
material-interface role or their exterior tag for boundary conditions.

The final audit independently checks the requested angle and area bounds,
constraint preservation, coverage, topology and constrained-Delaunay legality.
Successful refinement therefore meets the requested criteria, within the
documented floating-point audit margins ($10^{-8}$ degrees and $10^{-10}$
relative area). It does not estimate the magnetic field error. Check field
convergence with spacing and exterior-domain changes as before.

\subsection{Limits and failure diagnostics}
This is an original, bounded Lua implementation of segment-first Delaunay
refinement inspired by the algorithms studied in Triangle. It is not a full
port of Triangle's refinement machinery. In particular, it has no specialised
protection for acute input corners, no coarsening and no guaranteed smooth
gradation between local sizes. Candidate disk checks are conservative across
all segments, which can request more subdivisions than a visibility-aware test.

A material wedge with a smaller angle than the requested target cannot be
triangulated while satisfying that target everywhere. Closely spaced features
can require many extra points. Splitting stops with an error when the Steiner
budget, the total vertex cap or the geometric tolerance prevents further
progress; a candidate indistinguishable from an existing vertex also raises
an error. No partial result is silently accepted or saved as a successful mesh.
Reduce the target or review the geometry first; increase the budget only when
the extra resolution is intentional. These controls are shared by all three
TeX formats and use no external meshing program.
